\documentclass[../main]{subfiles}
\ifSubfilesClassLoaded{\setcounter{chapter}{1}}{}
\begin{document}

\chapter{Landasan Teori dan Konsep Dasar Kriptografi}

\begin{subcpmk}
  \item Sub-CPMK 1.1: Menjelaskan model komunikasi aman dan ancaman keamanan informasi
  \item Sub-CPMK 1.2: Mendefinisikan confidentiality, integrity, dan authenticity
  \item Sub-CPMK 1.3: Membedakan kriptografi simetris dan asimetris
\end{subcpmk}

% ============================================================
% MATERI POKOK
% ============================================================
\section{Pengantar Kriptografi}

\subsection{Definisi dan Sejarah}

Kriptografi berasal dari bahasa Yunani: \textit{kryptos} (tersembunyi) dan \textit{graphia} (menulis). Secara umum, \crypt{kriptografi} adalah ilmu dan seni untuk mengamankan informasi dengan mengubah pesan yang dapat dibaca (\crypt{plaintext}) menjadi bentuk yang tidak dapat dibaca tanpa kunci (\crypt{ciphertext}) \cite{schneier1996}.

Sejarah kriptografi dimulai sejak zaman kuno:
\begin{itemize}
  \item \textbf{Antik}: Scytale (Sparta), cipher substitusi sederhana
  \item \textbf{Abad Pertengahan}: Cipher Vigenère (abad ke-16)
  \item \textbf{Perang Dunia II}: Mesin Enigma (Jerman), pemecahan oleh Inggris
  \item \textbf{Era Modern}: DES (1977), RSA (1977), AES (2001)
\end{itemize}

Bukti earliest penggunaan kriptografi ditemukan dalam hieroglif non-standar yang diukir di dinding makam dari Kerajaan Lama Mesir sekitar 1900 SM. Menurut Wikipedia, tablet tanah liat dari Mesopotamia sekitar 1500 SM jelas dimaksudkan untuk melindungi informasi—satu tablet ditemukan mengenkripsi resep pematung untuk glasir keramik yang nilainya secara komersial. Para sarjana Ibrani menggunakan cipher substitusi monoalfabetik sederhana seperti Atbash cipher sekitar 600-500 SM, menunjukkan bahwa kebutuhan akan komunikasi rahasia telah ada sejak ribuan tahun lalu.

Evolusi kriptografi modern dimulai setelah Perang Dunia II dengan kemajuan signifikan dalam studi kriptografi, terutama pengenalan cipher kunci asimetris (kadang disebut cipher kunci publik). Menurut Wikipedia, cipher ini menggunakan dua kunci yang terkait secara matematis untuk enkripsi pesan yang sama, dan beberapa algoritma mengizinkan publikasi salah satu kunci karena sangat sulit menentukan satu kunci hanya dari pengetahuan kunci lainnya. Perkembangan internet untuk tujuan komersial pada tahun 1990-an menuntut standar enkripsi yang luas, yang pada akhirnya mengarah pada penggantian DES dengan AES melalui kompetisi publik yang diorganisir oleh NIST.

\subsection{Terminologi Dasar}

\begin{table}[htbp]
\centering
\begin{tabular}{|l|>{\raggedright\arraybackslash}p{8cm}|}
\hline
\textbf{Istilah} & \textbf{Definisi} \\
\hline
Plaintext & Pesan asli sebelum enkripsi \\
\hline
Ciphertext & Pesan terenkripsi, tidak dapat dibaca tanpa kunci \\
\hline
Cipher & Algoritma atau skema untuk enkripsi dan dekripsi \\
\hline
Key (Kunci) & Nilai rahasia yang mengontrol operasi enkripsi/dekripsi \\
\hline
Enkripsi & Proses mengubah plaintext menjadi ciphertext \\
\hline
Dekripsi & Proses mengubah ciphertext kembali ke plaintext \\
\hline
Kriptanalisis & Studi untuk memecahkan cipher tanpa mengetahui kunci \\
\hline
\end{tabular}
\caption{Terminologi Dasar Kriptografi}
\end{table}

\begin{konsep}
Dalam model dasar kriptografi, Alice mengirim pesan kepada Bob. Eve (eavesdropper) mencoba mencuri dengar. Enkripsi memastikan Eve tidak dapat memahami isi pesan meskipun berhasil menangkap ciphertext.
\end{konsep}

\section{Model Komunikasi Aman dan Ancaman}

\subsection{Model Komunikasi Kriptografi}

Model standar melibatkan:
\begin{itemize}
  \item \textbf{Alice (Pengirim)}: Mengirim plaintext, melakukan enkripsi
  \item \textbf{Bob (Penerima)}: Menerima ciphertext, melakukan dekripsi
  \item \textbf{Eve (Penyerang)}: Mencoba mencuri dengar (eavesdropping) atau memodifikasi pesan
  \item \textbf{Kanal}: Saluran komunikasi (internet, udara, jaringan) yang mungkin tidak aman
\end{itemize}

Tujuan kriptografi adalah memastikan komunikasi aman meskipun kanal tidak terpercaya. Model komunikasi ini mencerminkan skenario dunia nyata di mana informasi sensitif harus ditransmisikan melalui jaringan publik yang rentan terhadap intervensi pihak ketiga. Dalam implementasi praktis, Alice dan Bob dapat merepresentasikan berbagai entitas seperti pengguna dan server, dua sistem dalam jaringan korporat, atau bahkan proses yang berkomunikasi dalam satu sistem komputer.

\subsection{Ancaman Keamanan}

\begin{table}[htbp]
\centering
\small
\begin{tabular}{|l|p{6cm}|}
\hline
\textbf{Ancaman} & \textbf{Deskripsi} \\
\hline
Interception & Eve mencuri dengar pesan tanpa diketahui \\
\hline
Modification & Eve mengubah isi pesan dalam perjalanan \\
\hline
Fabrication & Eve membuat pesan palsu seolah dari Alice \\
\hline
Replay & Eve merekam pesan dan mengirim ulang di waktu lain \\
\hline
Masquerade & Eve berpura-pura sebagai pihak lain \\
\hline
\end{tabular}
\caption{Jenis Ancaman Keamanan}
\end{table}

Ancaman keamanan diklasifikasikan berdasarkan dampaknya terhadap layanan keamanan informasi \cite{rfc4949}. Ancaman terhadap \textbf{confidentiality} mencakup serangan langsung untuk mendapatkan akses ke sistem atau teknik man-in-the-middle di mana penyerang memposisikan diri dalam aliran informasi untuk mencegat data. Ancaman terhadap \textbf{integrity} berupa modifikasi tidak sah terhadap data atau log sistem untuk menyembunyikan serangan.

{\sloppy Untuk mengatasi ancaman tersebut, kriptografi menyediakan mekanisme: \textbf{confidentiality} (kerahasiaan), \textbf{integrity} (keutuhan), \textbf{availability} (ketersediaan), dan \textbf{non-repudiation} (ketidakbantahan).\par} Ketiga layanan pertama dikenal sebagai CIA triad; non-repudiation melengkapi model untuk transaksi yang memerlukan bukti hukum. Implementasi yang tepat memastikan komunikasi digital tetap aman dan terpercaya.

\section{Urgensi Kriptografi dalam Teknologi Informasi}

\subsection{Mengapa Komunikasi Publik Rentan?}

Hampir semua komunikasi modern---surat elektronik, SMS, browser web, perbankan daring---melewati kanal publik (internet, jaringan seluler, infrastruktur bersama). Tanpa proteksi kriptografi, pesan dapat disadap, diubah, atau dipalsukan. Karena itu kriptografi menjadi inti banyak teknologi keamanan siber: memahami perannya penting untuk memahami cara mengamankan ruang siber \cite{munir_slides_itb,stallings2017}.

\subsection{Kasus Ilustratif: Kebocoran dan Serangan}

Beberapa kasus berikut bersifat ilustratif (bukan laporan forensik lengkap). Tujuannya menunjukkan bahwa kegagalan melindungi data berdampak nyata pada privasi, kepercayaan publik, dan ekonomi digital \cite{munir_slides_itb}:
\begin{itemize}
  \item \textbf{Kebocoran data pribadi berulang}: data peserta asuransi kesehatan, pengguna e-commerce, atau nasabah lembaga keuangan yang bocor ke publik menekankan pentingnya enkripsi, kontrol akses, dan tata kelola data.
  \item \textbf{Penyadapan atau akses tidak sah terhadap fasilitas komunikasi}: menggarisbawahi kebutuhan confidentiality pada kanal diplomatik maupun korporat.
  \item \textbf{Peretasan situs lembaga atau ransomware}: dokumen rahasia atau data nasabah berisiko diekspos atau dienkripsi paksa oleh penyerang---menunjukkan bahwa keamanan tidak hanya soal firewall, tetapi juga proteksi kriptografi dan manajemen kunci.
\end{itemize}

\begin{konsep}
Pelajaran utama: tanpa kerahasiaan, integritas, dan otentikasi yang memadai, data di kanal publik maupun di penyimpanan mudah menjadi sasaran. Kriptografi bukan satu-satunya solusi keamanan, tetapi merupakan fondasi teknis yang sulit digantikan.
\end{konsep}

\subsection{Enkripsi At Rest dan In Transit}

Dua aplikasi utama enkripsi dalam praktik TI \cite{munir_slides_itb}:
\begin{itemize}
  \item \textbf{Encryption at rest}: melindungi data yang disimpan (disk, basis data, berkas dokumen) agar tetap rahasia jika media fisik atau salinan cadangan dicuri.
  \item \textbf{Encryption in transit} (\textit{on motion}): melindungi pesan
        yang dikirim melalui jaringan.
        Contoh familiar adalah enkripsi ujung-ke-ujung pada aplikasi pesan
        atau HTTPS/TLS pada browser.
\end{itemize}

Gambar~\ref{fig:at-rest-in-transit} memberi contoh keduanya dari aplikasi sehari-hari.

\begin{figure}[htbp]
\centering
\begin{minipage}[b]{0.42\textwidth}
  \centering
  \includegraphics[width=\linewidth]{bab-02/encrypt-document}\\[2pt]
  {\small (a) At rest: enkripsi berkas dokumen dengan kata sandi}
\end{minipage}\hfill
\begin{minipage}[b]{0.48\textwidth}
  \centering
  \includegraphics[width=\linewidth]{bab-02/whatsapp-e2e}\\[2pt]
  {\small (b) In transit: enkripsi ujung-ke-ujung pada aplikasi pesan}
\end{minipage}
\caption{Contoh encryption at rest dan encryption in transit \cite{munir_slides_itb}.}
\label{fig:at-rest-in-transit}
\end{figure}

\begin{catatan}
Sistem modern sering menggabungkan keduanya: data dienkripsi saat disimpan \textit{dan} saat dikirim. TLS (Bab~VIII) adalah contoh proteksi in transit; mode penyimpanan seperti XTS (Bab~V) berkaitan dengan proteksi at rest.
\end{catatan}

\subsection{Lembaga Terkait Kriptografi di Indonesia}

Di Indonesia, pengembangan kapasitas sandi dan keamanan siber melibatkan lembaga negara dan pendidikan antara lain \cite{munir_slides_itb,bssn_official}:
\begin{enumerate}
  \item \textbf{Badan Siber dan Sandi Negara (BSSN)} --- \url{https://bssn.go.id} --- lembaga yang menggabungkan fungsi sandi negara dan aspek aplikasi informatika terkait keamanan siber.
  \item \textbf{Politeknik Siber dan Sandi Negara (PSSN)} --- \url{https://poltekssn.ac.id} --- pendidikan tinggi kedinasan di bidang siber dan sandi.
\end{enumerate}

Museum Sandi di Yogyakarta (\url{https://museum.lemsaneg.go.id/}, Jl.~Faridan Muridan Noto No.~21, Kota Baru) menyimpan koleksi alat sandi yang pernah digunakan di Indonesia dan menjadi referensi sejarah praktis bagi mahasiswa yang ingin melihat artefak kriptografi klasik secara langsung \cite{munir_slides_itb}. Salah satu koleksinya adalah mesin sandi SN-101 (Gambar~\ref{fig:museum-sandi}).

\begin{figure}[htbp]
\centering
\begin{minipage}[b]{0.48\textwidth}
  \centering
  \includegraphics[width=\linewidth]{bab-02/museum-sandi}\\[2pt]
  {\small (a) Gedung Museum Sandi, Yogyakarta}
\end{minipage}\hfill
\begin{minipage}[b]{0.44\textwidth}
  \centering
  \includegraphics[width=\linewidth]{bab-02/mesin-sandi-sn101}\\[2pt]
  {\small (b) Mesin sandi SN-101}
\end{minipage}
\caption{Museum Sandi dan salah satu koleksi alat sandinya \cite{munir_slides_itb}.}
\label{fig:museum-sandi}
\end{figure}

Bagian berikutnya merinci \textbf{layanan keamanan informasi} yang menjawab empat masalah Alice--Bob dan ancaman yang telah dibahas.

\section{Layanan Keamanan Informasi}

Empat masalah Alice--Bob (penyadapan, penyamaran, manipulasi, penyangkalan) dijawab oleh layanan keamanan yang didukung kriptografi \cite{munir_slides_itb,rfc4949,nist_cia}: \textbf{confidentiality}, \textbf{authentication}/\textbf{authenticity}, \textbf{integrity}, dan \textbf{non-repudiation}, dilengkapi \textbf{availability} dalam model CIA. Kasus kebocoran data dan serangan pada bagian urgensi memperkuat motivasi bahwa layanan ini bukan sekadar istilah teoretis, melainkan kebutuhan praktis. Tiga layanan pertama CIA triad dikenal luas; non-repudiation dan authenticity melengkapi model untuk transaksi yang memerlukan kepercayaan dan bukti hukum.

\subsection{Confidentiality (Kerahasiaan)}

Memastikan informasi hanya dapat diakses oleh pihak yang berwenang. Plaintext tidak boleh terbaca oleh Eve meskipun dia memiliki ciphertext. \textbf{Cara mencapai}: enkripsi dengan kunci yang hanya diketahui Alice dan Bob \cite{stallings2017}. Gambar~\ref{fig:layanan-kerahasiaan} menunjukkan plainteks yang menjadi tidak bermakna setelah dienkripsi.

\begin{figure}[htbp]
\centering
\includegraphics[width=0.8\textwidth]{bab-02/kerahasiaan-plain-cipher}
\caption{Kerahasiaan: plainteks (a) dienkripsi menjadi cipherteks (b) yang tidak dapat dipahami pihak yang tidak berhak \cite{munir_slides_itb}.}
\label{fig:layanan-kerahasiaan}
\end{figure}

Confidentiality merupakan fondasi keamanan informasi. Akses ke informasi harus dikontrol untuk mencegah pembagian data tidak sah. Implementasi efektif memerlukan enkripsi data, kontrol akses, dan autentikasi. Ancaman terhadap confidentiality mencakup serangan siber maupun kesalahan manusia seperti berbagi kredensial atau kegagalan mengenkripsi komunikasi.

\subsection{Integrity (Keutuhan)}

Memastikan informasi tidak diubah selama transmisi. Bob harus yakin pesan yang diterima sama persis dengan yang dikirim Alice. \textbf{Cara mencapai}: fungsi hash, Message Authentication Code (MAC), atau tanda tangan digital. Perubahan kecil pun harus terdeteksi, seperti satu digit suhu atau satu bagian foto pada Gambar~\ref{fig:layanan-integritas}.

\begin{figure}[htbp]
\centering
\includegraphics[width=\textwidth]{bab-02/integritas-pesan-diubah}
\caption{Integritas: pesan teks dan citra yang diubah sedikit tampak hampir sama dengan aslinya \cite{munir_slides_itb}.}
\label{fig:layanan-integritas}
\end{figure}

Integrity memastikan data tetap tepercaya dan bebas dari perusakan. Setiap perubahan tidak sah terhadap data harus terdeteksi dan dicegah. Pelanggaran integrity dapat menyebabkan keputusan bisnis yang salah atau kehilangan kepercayaan. Teknologi hash, tanda tangan digital, dan MAC menyediakan mekanisme verifikasi yang diperlukan.

\subsection{Availability (Ketersediaan)}

Memastikan sistem dan data dapat diakses oleh pihak yang berwenang ketika diperlukan. Serangan denial-of-service (DoS) mengancam availability dengan membuat layanan tidak dapat diakses.

Kriptografi secara tidak langsung mendukung availability melalui proteksi terhadap serangan yang menargetkan infrastruktur. Manajemen kunci dan redundansi sistem juga berperan menjaga ketersediaan layanan kriptografi. Dalam praktik, availability sering ditangani melalui kombinasi kebijakan operasional, redundansi jaringan, dan mitigasi DoS.

\subsection{Non-Repudiation (Ketidakbantahan)}

Memastikan pihak yang melakukan tindakan tidak dapat menyangkal bahwa ia telah melakukannya. Alice tidak dapat menyangkal telah mengirim pesan; Bob tidak dapat menyangkal telah menerimanya. \textbf{Cara mencapai}: tanda tangan digital, sertifikat digital, log audit. Gambar~\ref{fig:layanan-nirpenyangkalan} mengilustrasikan nasabah yang menyangkal instruksi transfernya sendiri.

\begin{figure}[htbp]
\centering
\includegraphics[width=0.6\textwidth]{bab-02/nirpenyangkalan-transfer}
\caption{Non-repudiation: tanpa bukti kriptografis, A dapat menyangkal telah meminta transfer ke B \cite{munir_slides_itb}.}
\label{fig:layanan-nirpenyangkalan}
\end{figure}

Non-repudiation penting untuk transaksi hukum, kontrak elektronik, dan pembayaran daring. Tanda tangan digital memberikan bukti matematis yang tidak dapat dipalsukan tentang asal dan integritas pesan. Sertifikat digital dari CA tepercaya mendukung verifikasi identitas dalam skala global \cite{stallings2017}.

\subsection{Authenticity (Keaslian)}

Authenticity memastikan pengirim pesan benar-benar pihak yang mengklaim. Bob yakin pesan berasal dari Alice, bukan Eve yang berpura-pura. \textbf{Cara mencapai}: tanda tangan digital, sertifikat digital, protokol autentikasi.

\begin{figure}[!htb]
\centering
\includegraphics[width=0.5\textwidth]{bab-02/autentikasi-penyamaran}
\caption{Autentikasi: pihak ketiga di tengah mengirim sinyal palsu seolah-olah berasal dari A, sehingga B perlu memverifikasi asal pesan \cite{munir_slides_itb}.}
\label{fig:layanan-autentikasi}
\end{figure}

Gambar~\ref{fig:layanan-autentikasi} menggambarkan masalah ini secara sederhana: B menerima sinyal asap, tetapi tidak dapat memastikan apakah sinyal itu dari A atau dari penyamar.
\begin{table}[htbp]
\centering
\small
\begin{tabular}{|l|>{\raggedright\arraybackslash}p{4.8cm}|>{\raggedright\arraybackslash}p{4.5cm}|}
\toprule
\textbf{Layanan} & \textbf{Definisi} & \textbf{Mekanisme Kriptografi} \\
\midrule
Confidentiality & Hanya pihak berwenang mengakses & Enkripsi \\
Integrity & Data tidak diubah tanpa izin & Hash, MAC \\
Availability & Data/sistem dapat diakses saat dibutuhkan & Manajemen kunci, redundansi \\
Non-repudiation & Tindakan tidak dapat disangkal & Tanda tangan digital \\
Authenticity & Identitas pengirim/penerima terverifikasi & Sertifikat, protokol auth \\
\bottomrule
\end{tabular}
\caption{Layanan keamanan informasi dan mekanisme kriptografi.}
\end{table}

\begin{contoh}
Dalam pembayaran daring: confidentiality melindungi data kartu kredit; integrity memastikan jumlah transaksi tidak diubah; authenticity memastikan pembayaran dari pemilik kartu yang sah; non-repudiation mencegah penyangkalan transaksi.
\end{contoh}

\section{Kriptologi: Kriptografi dan Kriptanalisis}

Banyak orang menganggap bahwa kriptografi hanyalah tentang membuat kode rahasia. Namun, dalam perspektif akademis, kriptografi adalah bagian dari bidang yang lebih luas yang disebut \textbf{Kriptologi}.

\subsection{Hubungan antara Kriptografi dan Kriptanalisis}

Kriptologi adalah studi ilmiah mengenai keamanan informasi yang terbagi menjadi dua cabang utama yang saling berlawanan namun saling melengkapi:

\begin{enumerate}
    \item \textbf{Kriptografi (\textit{Cryptography})}: Ilmu dan seni untuk membangun sistem penyandian yang aman. Fokus utamanya adalah menciptakan algoritma yang mampu menjaga kerahasiaan, integritas, dan otentikasi pesan.
    \item \textbf{Kriptanalisis (\textit{Cryptanalysis})}: Ilmu dan seni untuk memecahkan cipherteks menjadi plainteks tanpa mengetahui kunci yang digunakan. Fokus utamanya adalah mencari celah, kelemahan, atau pola dalam algoritma untuk membobol keamanan.
\end{enumerate}

Hubungan antara keduanya adalah sebuah siklus evolusi. Ketika seorang kriptografer menciptakan cipher baru yang dianggap aman, para kriptanalis akan mencoba memecahkannya. Jika berhasil, maka cipher tersebut dinyatakan tidak aman, dan kriptografer harus menciptakan metode baru yang lebih kuat. Itulah sebabnya kriptanalisis sering disebut sebagai "lawan" sekaligus "guru" bagi kriptografi.

\subsection{Sejarah Singkat Kriptanalisis}

Kriptanalisis sudah ada sejak abad ke-9, jauh sebelum komputer ditemukan. Salah satu tokoh paling berpengaruh adalah \textbf{Al-Kindi}, seorang ilmuwan Arab yang pertama kali memperkenalkan teknik \textbf{analisis frekuensi}. Ia menyadari bahwa dalam setiap bahasa, huruf-huruf tertentu muncul lebih sering daripada huruf lainnya. Dengan menghitung frekuensi kemunculan karakter dalam cipherteks, Al-Kindi dapat menebak huruf apa yang disandikan.

Sejarah mencatat bahwa keberhasilan kriptanalisis seringkali mengubah jalannya sejarah dunia. Contoh paling nyata adalah pemecahan kode \textit{Enigma} oleh pihak Sekutu pada Perang Dunia II, yang diperkirakan telah memperpendek durasi perang dan menyelamatkan jutaan nyawa.

\subsection{Kriptografi di Indonesia}

Di Indonesia, pengelolaan keamanan informasi dan sandi dikelola oleh lembaga negara yang kompeten, yaitu \textbf{Badan Siber dan Sandi Negara (BSSN)}. BSSN bertanggung jawab atas koordinasi keamanan siber dan kriptografi nasional. Selain itu, terdapat \textbf{Politeknik Siber dan Sandi Negara (PSSN)} sebagai lembaga pendidikan khusus yang mencetak tenaga ahli di bidang sandi dan keamanan siber.

Bagi mahasiswa yang tertarik mendalami sejarah kriptografi secara fisik, Indonesia memiliki \textbf{Museum Sandi} di Yogyakarta. Museum ini merupakan satu-satunya museum sandi di dunia yang menyimpan berbagai koleksi alat sandi bersejarah, mulai dari alat manual hingga mesin sandi mekanik yang pernah digunakan dalam perjuangan kemerdekaan dan pemerintahan Indonesia.

\begin{figure}[htbp]
  \centering
  \includegraphics[width=0.6\textwidth]{referensi/01-Pengantar-Kriptografi-2026/images/llm/page_065_img_01.png}
  \caption{Museum Sandi Yogyakarta: Pusat pelestarian sejarah kriptografi di Indonesia.}
\end{figure}


% ============================================================
% AKTIVITAS PEMBELAJARAN
% ============================================================

\begin{aktivitas}
  \item \textbf{Diskusi Kelompok}: Identifikasi 5 aplikasi sehari-hari yang menggunakan kriptografi (misalnya: WhatsApp, perbankan online, e-commerce). Jelaskan ancaman apa yang dilindungi masing-masing.
  
  \item \textbf{Analisis Kasus}: Diberikan skenario: Eve mencuri ciphertext dari komunikasi Alice dan Bob. Apa yang bisa dan tidak bisa dilakukan Eve? Jelaskan peran confidentiality, integrity, dan authenticity.
  
  \item \textbf{Perbandingan}: Buat tabel perbandingan antara kriptografi simetris dan asimetris untuk penggunaannya dalam: enkripsi file, pertukaran kunci, tanda tangan digital.
  
  \item \textbf{Mind Mapping}: Buat mind map yang menghubungkan confidentiality, integrity, authenticity dengan mekanisme kriptografi yang mendukungnya.
\end{aktivitas}

% ============================================================
% LATIHAN DAN REFLEKSI
% ============================================================

\begin{latihan}
  \item Jelaskan definisi plaintext, ciphertext, cipher, dan key! Berikan contoh konkret.
  
  \item Sebutkan dan jelaskan tiga pilar keamanan informasi (confidentiality, integrity, authenticity)! Berikan contoh ancaman yang melanggar masing-masing.
  
  \item Apa perbedaan utama antara kriptografi simetris dan asimetris? Kapan masing-masing lebih cocok digunakan?
  
  \item Jelaskan model komunikasi kriptografi (Alice, Bob, Eve). Apa peran masing-masing?
  
  \item Berikan 3 contoh ancaman keamanan dan jelaskan bagaimana kriptografi dapat mengatasinya.
  
  \item \textbf{Refleksi}: Bagaimana pemahaman Anda tentang keamanan informasi sebelum dan sesudah mempelajari bab ini?
\end{latihan}

% ============================================================
% ASESMEN
% ============================================================

\begin{asesmen}
\textbf{Instrumen Penilaian untuk Sub-CPMK 1.1, 1.2, 1.3}

\textbf{A. Pilihan Ganda}

\begin{enumerate}
  \item Manakah yang BUKAN merupakan pilar keamanan informasi?
  \begin{enumerate}
    \item Confidentiality
    \item Integrity
    \item Complexity
    \item Authenticity
  \end{enumerate}
  
  \item Dalam kriptografi simetris, jumlah kunci yang digunakan adalah:
  \begin{enumerate}
    \item Satu kunci untuk enkripsi dan dekripsi
    \item Dua kunci berbeda
    \item Tidak ada kunci
    \item Kunci publik saja
  \end{enumerate}
  
  \item Fungsi utama enkripsi adalah untuk mencapai:
  \begin{enumerate}
    \item Integritas
    \item Kerahasiaan (confidentiality)
    \item Keaslian
    \item Ketersediaan
  \end{enumerate}
\end{enumerate}

\textbf{B. Essay}

\begin{enumerate}
  \item Jelaskan dengan kata-kata Anda sendiri apa yang dimaksud dengan kriptografi dan berikan 2 contoh penggunaannya dalam kehidupan nyata.
  
  \item Bandingkan kriptografi simetris dan asimetris. Mana yang lebih cocok untuk enkripsi file besar dan mengapa?
\end{enumerate}

\textbf{Rubrik Penilaian}: Lihat Lampiran A
\end{asesmen}

% ============================================================
% CHECKLIST KOMPETENSI
% ============================================================

\begin{checklist}
  \item Saya dapat menjelaskan terminologi dasar kriptografi (plaintext, ciphertext, cipher, key)
  \item Saya dapat menjelaskan model komunikasi aman dan ancaman keamanan
  \item Saya dapat mendefinisikan confidentiality, integrity, dan authenticity
  \item Saya dapat membedakan kriptografi simetris dan asimetris
  \item Saya dapat memberikan contoh penggunaan masing-masing jenis kriptografi
\end{checklist}

% ============================================================
% RANGKUMAN
% ============================================================

\begin{rangkuman}
Bab ini membahas landasan teori kriptografi, termasuk definisi, sejarah, terminologi, model komunikasi aman, ancaman keamanan, tiga pilar keamanan informasi, serta perbedaan kriptografi simetris dan asimetris.

\textbf{Poin Kunci:}
\begin{itemize}
  \item Kriptografi mengamankan informasi melalui enkripsi dan dekripsi
  \item Tiga pilar: Confidentiality, Integrity, Authenticity
  \item Simetris: satu kunci, cepat; Asimetris: pasangan kunci, handal untuk distribusi
  \item Pemahaman landasan teori adalah fondasi untuk mempelajari algoritma kriptografi
\end{itemize}

\textbf{Kata Kunci}: \crypt{Kriptografi}, \crypt{Plaintext}, \crypt{Ciphertext}, \crypt{Enkripsi}, \crypt{Dekripsi}, \crypt{Confidentiality}, \crypt{Integrity}, \crypt{Authenticity}
\end{rangkuman}

\ifSubfilesClassLoaded{
  \renewcommand{\bibname}{Daftar Pustaka}
  \bibliographystyle{plain}
  \bibliography{../references}
}{}
\end{document}
